% !TeX root = ../main.tex
\section{Methodology}\label{sec:method}
\rubric{Experimental rigor}{25}{Pinning and isolation, repetitions, documented settings and reproducible scripts.}

\subsection{Platform}
\newcommand{\acspinning}{single-thread runs and the loaded-latency probe pinned to CPU~2 (a physical
core, not CPU~0); multi-threaded runs place one thread per physical core}
\newcommand{\acsenvextra}{%
  Memory configuration & \sysMemory{} as reported by the kernel (\code{MemTotal}); single NUMA node;
  DIMM count and speed not probed (needs root/\code{dmidecode}) \\
}
All measurements were taken on the desktop system in \cref{tab:environment}: a ten-core Intel
i9-10900F (Comet Lake, Skylake-family core) with per-core \qty{32}{\kibi\byte} L1d and
\qty{256}{\kibi\byte} L2, a shared \qty{20}{\mebi\byte} L3, and DDR4 memory on one NUMA node. The table is
generated by \code{scripts/collect_env.sh} on the test host; each raw-data directory also stores its own
snapshot (\code{env.txt}) taken immediately before that experiment, so the settings of every figure are on
record even if the machine changes later.
%% Measurement environment table (common requirement: measurement discipline).
%% Values come from common/system_info.tex. Before \input-ing this file a project may define
%%   \acspinning   how runs were pinned/isolated
%%   \acsenvextra  extra rows, e.g.  \newcommand{\acsenvextra}{SSD & model, firmware, PCIe link \\}
\providecommand{\acspinning}{\todo{e.g.\ \code{taskset -c 2} on a P-core; isolcpus if used}}
\providecommand{\acsenvextra}{}
\begin{table}[tbp]
  \centering\small
  \caption{Measurement environment.}\label{tab:environment}
  \begin{tabularx}{\linewidth}{@{}>{\raggedright\arraybackslash}p{0.2\linewidth}X@{}}
    \toprule
    CPU               & \sysCPU{} (\sysArch) \\
    SIMD ISA          & \sysISA \\
    Topology          & \sysCores; NUMA: \sysNUMA \\
    SMT               & \sysSMT \\
    Caches            & \sysCaches \\
    Memory            & \sysMemory; DIMMs: \sysDIMM \\
    Pages             & base \sysPageSize; THP: \sysHugePages \\
    OS / kernel       & \sysOS; \sysKernel \\
    Frequency policy  & governor: \sysGovernor; driver: \sysDriver; turbo: \sysTurbo; EPP: \sysEPP \\
    Pinning / isolation & \acspinning; \sysIsolation \\
    Compilers         & \sysGCC; \sysClang \\
    Tools             & \sysPerf{} (perf\_event\_paranoid \sysPerfParanoid); \sysFio; MLC: \sysMLC \\
    GPU (optional)    & \sysGPU; CUDA \sysCUDA \\
    \acsenvextra
    Collected         & \sysCollected \\
    \bottomrule
  \end{tabularx}
\end{table}


\subsection{Tools}
All benchmarks share the harness in \code{ACS-Code/common} (monotonic raw timer, CPU pinning) plus a
project header \code{src/memprof.h} (page-policy allocation, ring generators, counters), and raw
per-repetition rows are reduced by \code{tools/summarize} (median and inter-quartile range).
\code{scripts/derive.py} then computes every derived number quoted in the text (plateaus, knee, Little's
Law, AMAT fit) into \code{summary.csv}; the report reads those values directly, so no number in the prose is
copied by hand.
\begin{knobs}
  \code{mem_latency}   & single-thread pointer-chase latency: working-set size, access size, node stride, page size, ring generator, timer, counters \\
  \code{mem_bandwidth} & streaming bandwidth with $T$ pinned threads: read/write/copy/triad/rwmix, access-size blocks, byte stride \\
  \code{mem_loaded}    & loaded latency: $T$ load-generator cores (stream, rwmix, or $k$ independent chases) plus a concurrent latency probe \\
  \code{kernel_miss}   & SAXPY loop whose footprint, order and page size vary, with counters \\
  Counters             & \code{perf_event_open(2)} called directly, user space only \cite{linux_perf} \\
  Intel MLC            & not installed; substitute cross-checks in \cref{sec:validation} \cite{intel_mlc} \\
\end{knobs}
The \code{perf} command-line tool is not installed and \code{perf_event_paranoid} is
\sysPerfParanoid{}, which still allows a process to count its \emph{own} user-mode events. The benchmarks
therefore open the counters themselves (\code{exclude_kernel=1}, \code{exclude_hv=1}): cycles,
instructions, L1D read misses, LLC read references and misses, and dTLB read misses (completed page walks).
Events are opened independently so the kernel may multiplex them; counts are scaled by
\code{time_enabled/time_running} and the smallest running fraction is recorded per row
(\code{ctr_cover}). Intel MLC is the spec's suggested cross-check but is not installed and could not be
installed for this study; \code{run.sh mlc} records that fact in its raw directory, and
\cref{sec:validation} describes the clearly labelled substitutes.

\subsection{Controls}
\begin{itemize}
  \item \textbf{Pinning and isolation:} \acspinning. SMT: \sysSMT. The machine is otherwise idle and
        experiments are run one at a time.
  \item \textbf{Frequency:} governor \sysGovernor, turbo \sysTurbo. Cycle figures use the effective
        core frequency each run measures in-process from a dependent add chain
        (\valqty{zq_ghz}{\giga\hertz} median), not a nominal clock.
  \item \textbf{Pages:} every buffer is a \qty{2}{\mebi\byte}-aligned anonymous mapping whose page size is
        chosen by \code{madvise} \emph{before first touch}: \code{MADV_HUGEPAGE} for ``huge'' and
        \code{MADV_NOHUGEPAGE} for ``4k''. The second is essential because THP is \sysHugePages{}: without
        it a ``\qty{4}{\kibi\byte}'' baseline can silently receive huge pages. Each run records the
        \code{AnonHugePages} actually granted (from \code{/proc/self/smaps_rollup}).
  \item \textbf{Pre-faulting and warm-up:} every buffer is written once to fault in its pages, then the
        ring or stream is traversed (or one untimed repetition is run) before timing, so neither page
        faults nor cold-cache effects fall inside the timed region.
  \item \textbf{Repetitions:} 8 per point; figures show the median with inter-quartile-range error bars.
  \item \textbf{Timing:} each timed interval is calibrated to $\geq$\qtyrange{40}{200}{\milli\second},
        far above the timer resolution, so timer overhead is negligible. The loaded-latency windows read
        generator progress counters at the start and end of the \emph{same} interval in which the probe
        runs, so throughput and latency describe the same moment.
  \item \textbf{Provenance:} each \code{data/raw/<experiment>/<timestamp>/} directory holds the raw rows,
        \code{env.txt}, \code{build.txt} (git commit, compiler, flags, THP and counter settings),
        \code{commands.txt} (every command line) and \code{validate.txt} (the sanity gate of
        \cref{sec:validation}, run immediately before the experiment).
\end{itemize}
